\documentclass[a4paper]{article}

% Basic packages
% Español (México) con XeLaTeX: fontspec para fuentes Unicode, polyglossia
% para la división silábica y los nombres del documento en la variante mexicana.
\usepackage{fontspec}
\usepackage{polyglossia}
\setdefaultlanguage[variant=mexican]{spanish}
% Los nombres chinos van como «pinyin (original)»: xeCJK da a los caracteres
% Han su propia fuente; sin ella XeLaTeX los omite con un aviso.
% FandolSong no cubre algunos caracteres de nombres (U+73FA); WenQuanYi Zen Hei sí.
\usepackage[AutoFallBack=true]{xeCJK}
\setCJKmainfont{FandolSong-Regular.otf}
\setCJKfallbackfamilyfont{\CJKrmdefault}{WenQuanYi Zen Hei}
% polyglossia no traduce los nombres de listings.
\AtBeginDocument{\providecommand{\lstlistingname}{}\renewcommand{\lstlistingname}{Listado}%
  \providecommand{\lstlistlistingname}{}\renewcommand{\lstlistlistingname}{Índice de listados}}
\usepackage{amsmath, amssymb}
\usepackage{graphicx}
\usepackage[margin=2.5cm]{geometry}
\usepackage[most]{tcolorbox}
\usepackage{etoolbox}
\usepackage{listings}
\usepackage{hyperref}
\usepackage{booktabs}
\usepackage{subcaption}
\usepackage{float}

% Highlight boxes
\newtcolorbox{knowledgebox}[1]{
    enhanced, colback=blue!5!white, colframe=blue!75!black, colbacktitle=blue!75!black,
    coltitle=white, fonttitle=\bfseries, title={#1}, attach boxed title to top left={yshift=-2mm, xshift=2mm},
    boxrule=1pt, sharp corners
}

\newtcolorbox{importantbox}[1]{
    enhanced, colback=yellow!10!white, colframe=yellow!80!black, colbacktitle=yellow!80!black,
    coltitle=black, fonttitle=\bfseries, title={#1}, sharp corners
}

\newtcolorbox{warningbox}[1]{
    enhanced, colback=red!5!white, colframe=red!75!black, colbacktitle=red!75!black,
    coltitle=white, fonttitle=\bfseries, title={#1}, sharp corners
}

% Listings
\lstset{
    language=Python,
    basicstyle=\ttfamily\small,
    keywordstyle=\color{blue},
    stringstyle=\color{red!60!black},
    commentstyle=\color{green!60!black},
    breaklines=true,
    frame=single,
    numbers=left,
    numberstyle=\tiny\color{gray},
    captionpos=b,
    extendedchars=true
}

% Front-page metadata
\newcommand{\notetitle}{{[}LLM Agents F24{]} LLM Agents: Brief History and Overview — Shunyu Yao}
\newcommand{\noteauthors}{Elaboradas a partir de materiales públicos del curso}
\newcommand{\notedate}{16 de septiembre de 2024}
\newcommand{\videochannel}{Berkeley RDI}
\newcommand{\videopublishdate}{16 de septiembre de 2024}
\newcommand{\videoduration}{aproximadamente 75 minutos}
\newcommand{\videourl}{https://www.youtube.com/watch?v=RM6ZArd2nVc}
\newcommand{\videocoverpath}{cover.jpg}
\newcommand{\repourl}{https://github.com/hqhq1025/ai-course-notes}


% Footer with repo info
\usepackage{fancyhdr}
\pagestyle{fancy}
\fancyhf{}
\fancyfoot[C]{\thepage}
\fancyfoot[R]{\footnotesize\color{gray}\href{\repourl}{AI Course Notes}}
\renewcommand{\headrulewidth}{0pt}
\renewcommand{\footrulewidth}{0pt}

\begin{document}

\begin{titlepage}
\centering
{\Large Notas del curso\par}
\vspace{1.2cm}
{\huge\bfseries \notetitle\par}
\vspace{0.8cm}
{\large \noteauthors\par}
\vspace{0.3cm}
{\large \notedate\par}
\vspace{1.2cm}

\ifdefempty{\videocoverpath}{
{\small Al generar las notas, indique la ruta local de la portada del video.\par}
}{
\includegraphics[width=0.82\textwidth,height=0.45\textheight,keepaspectratio]{\videocoverpath}\par
}

\vfill
\begin{tcolorbox}[width=0.9\textwidth, colback=black!2!white, colframe=black!60, sharp corners]
\textbf{Autor o canal del video}: \videochannel\par
\textbf{Fecha de publicación}: \videopublishdate\par
\textbf{Duración del video}: \videoduration\par
\ifdefempty{\videourl}{}{
\textbf{Enlace del video}: \href{\videourl}{\nolinkurl{\videourl}}\par
}
\par
\textbf{Página del proyecto}: \href{\repourl}{\nolinkurl{\repourl}}\par
\end{tcolorbox}
\end{titlepage}

\tableofcontents
\newpage

%% --- Inicio del contenido --- %%

\section{Introducción: ¿qué es un LLM Agent?}

Shunyu Yao (姚顺雨) es el creador del marco ReAct y actualmente trabaja en OpenAI. Esta clase ofrece un recuento sistemático de la definición de los LLM Agent, la trayectoria histórica de su desarrollo y las direcciones futuras.

\subsection{La definición de Agent}

\begin{importantbox}{Las tres capas de la definición de Agent}
\begin{enumerate}
    \item \textbf{Text Agent}: un sistema inteligente que interactúa con el entorno, cuyas acciones (action) y observaciones (observation) se expresan en forma de lenguaje (text). El Text Agent no necesariamente usa un LLM; desde los inicios mismos de la IA ya existían Text Agent basados en reglas.
    \item \textbf{LM Agent}: un Text Agent que usa un modelo de lenguaje para generar sus acciones.
    \item \textbf{Reasoning Agent}: un Agent que usa un modelo de lenguaje para \textbf{razonar antes de actuar}, lo cual constituye un salto cualitativo.
\end{enumerate}
\end{importantbox}

Definir Agent requiere responder dos subpreguntas: ¿qué es la ``inteligencia'' (intelligence)? ¿Qué es el ``entorno'' (environment)?

\begin{itemize}
    \item El entorno puede ser físico (robots, conducción autónoma) o digital (videojuegos, páginas web, editores de código).
    \item La definición de ``inteligencia'' evoluciona con el tiempo: hace 60 años, un simple chatbot ya se consideraba inteligente, mientras que hoy ChatGPT ya no sorprende a nadie.
\end{itemize}

\subsection{Resumen de la sección}

La innovación central del LLM Agent no consiste en hacer que el modelo de lenguaje ejecute acciones, sino en hacer que \textbf{razone antes de actuar}. Lo que distingue al Reasoning Agent de todos los paradigmas de Agent anteriores es que el razonamiento (thinking) es una acción interna con un espacio de lenguaje ilimitado.
\section{Historia antes de los LLM Agent}

\subsection{Text Agent basados en reglas}

\begin{knowledgebox}{ELIZA (1966)}
Uno de los primeros chatbots, capaz de dar la impresión de «inteligencia» usando solo unas cuantas reglas (preguntar insistentemente o repetir lo que decía el usuario). Sin embargo, el método basado en reglas tiene limitaciones serias:
\begin{itemize}
    \item Altamente específico de dominio (domain-specific)
    \item Cada dominio nuevo requiere volver a escribir las reglas
    \item No se puede escalar a dominios abiertos complejos
\end{itemize}
\end{knowledgebox}

\subsection{Text Agent basados en aprendizaje por refuerzo}

Antes de la aparición de los LLM, otro paradigma popular era construir Text Agent con RL (por ejemplo, agentes para juegos de texto):

\begin{itemize}
    \item Tratar las observaciones y acciones de texto como los píxeles y las entradas de teclado en un videojuego
    \item Usar RL para optimizar la señal de recompensa
    \item \textbf{Limitación}: específico de dominio, requiere una señal de recompensa precisa, costo de entrenamiento alto
\end{itemize}

\subsection{El cambio revolucionario de los LLM}

Los LLM, entrenados únicamente con next-token prediction sobre enormes cantidades de texto, pueden resolver diversas tareas nuevas en el momento de la inferencia mediante prompting. Esta \textbf{generalidad} y esta \textbf{capacidad de aprendizaje con pocos ejemplos} los convierten en el «cerebro» ideal para construir Agent.

\subsection{Resumen de la sección}

Del symbolic AI agent basado en reglas al deep agent basado en RL, y de ahí al reasoning agent basado en LLM, la representación interna del Agent pasó de la lógica simbólica a los embeddings neuronales y de ahí al lenguaje natural, avanzando paso a paso hacia una mayor flexibilidad y generalidad.
\section{La fusión de razonamiento y acción: ReAct}

\subsection{El paradigma de Reasoning y el paradigma de Acting}

Antes de que apareciera ReAct, las aplicaciones de los LLM se dividían en dos paradigmas independientes:

\begin{itemize}
    \item \textbf{Reasoning} (como Chain-of-Thought): el LLM razona internamente paso a paso, lo que aumenta el test-time compute, pero no puede acceder a conocimiento externo ni a herramientas.
    \item \textbf{Acting} (como RAG, tool use): el LLM invoca el entorno externo (motores de búsqueda, calculadoras, API) para obtener conocimiento y capacidad de cómputo, pero carece de razonamiento para adaptarse a situaciones complejas.
\end{itemize}

Cada paradigma resuelve por separado un subproblema distinto de las preguntas y respuestas (problemas de razonamiento frente a problemas de conocimiento frente a problemas de cómputo), lo que produce métodos fragmentados (piecemeal).

\subsection{ReAct: un marco unificado de razonamiento y acción}

\begin{importantbox}{La idea central de ReAct}
Hacer que el LLM alterne entre generar un \textbf{pensamiento} (thought) y una \textbf{acción} (action), mientras el entorno devuelve una observación (observation). Basta con escribir un ejemplo de trayectoria como prompt que muestre cómo pensar y actuar para resolver la tarea, y el LLM puede imitar ese patrón.
\end{importantbox}

Flujo de trabajo concreto:
\begin{enumerate}
    \item El LLM genera un thought (por ejemplo, «necesito buscar la capitalización de mercado de estas tres empresas»)
    \item El LLM genera una action (por ejemplo, \texttt{Google{[}Apple market cap{]}})
    \item El entorno devuelve una observation (el resultado de la búsqueda)
    \item El thought, la action y la observation se agregan al contexto y el LLM continúa generando el siguiente paso
    \item Se repite hasta completar la tarea
\end{enumerate}

\begin{knowledgebox}{El beneficio mutuo entre razonamiento y acción}
\begin{itemize}
    \item \textbf{La acción ayuda al razonamiento}: obtiene información en tiempo real del exterior y ejecuta cálculos
    \item \textbf{El razonamiento ayuda a la acción}: planea estrategias según la situación actual y reajusta el plan ante anomalías. Por ejemplo, cuando la búsqueda devuelve el precio de la acción en lugar de la capitalización de mercado, el razonamiento puede indicar que el siguiente paso es buscar el número de acciones para calcularla de forma indirecta.
\end{itemize}
\end{knowledgebox}

\subsection{La generalidad de ReAct}

ReAct no se limita a las preguntas y respuestas: cualquier tarea que pueda convertirse en un «juego de texto» puede usar el marco de ReAct:

\begin{itemize}
    \item Navegación web, compras en línea
    \item Videojuegos (convirtiendo los píxeles en texto mediante image captioning)
    \item Simulación de tareas domésticas (ALFWorld)
    \item Ingeniería de software (SWE-bench)
\end{itemize}

\begin{warningbox}{Las limitaciones de un Acting Agent sin razonamiento}
En el juego ALFWorld, un agente puramente de acting (que no genera thought) repite la misma acción una y otra vez al encontrar una anomalía (por ejemplo, cuando el objeto buscado no está en la posición esperada), sin poder adaptarse. Al agregar la thinking action, el agente puede reflexionar sobre la situación actual y replanificar.
\end{warningbox}

\subsection{Resumen de la sección}

La aportación de ReAct no es solo un método concreto, sino la propuesta del paradigma del Reasoning Agent, que trata el razonamiento como una «acción interna» especial del agente. Este paradigma supera sistemáticamente tanto al razonamiento puro como a la acción pura.

\section{Significado teórico del Reasoning Agent}

\subsection{El razonamiento como acción interna de dimensión infinita}

El espacio de acciones de un Agent tradicional (ya sea symbolic o RL) está definido de manera fija por el entorno (por ejemplo, arriba, abajo, izquierda, derecha), mientras que el Reasoning Agent introduce un tipo de acción completamente nuevo: \textbf{el razonamiento}:

\begin{itemize}
    \item El razonamiento puede ser lenguaje natural de cualquier longitud; el espacio de acciones es \textbf{infinito}
    \item El razonamiento no cambia el entorno externo, solo cambia el contexto del propio Agent (su estado interno)
    \item El razonamiento influye en las decisiones de acción posteriores
\end{itemize}

\subsection{Comparación de los tres paradigmas de Agent}

\begin{table}[H]
\centering
\begin{tabular}{lccc}
\toprule
& \textbf{Symbolic AI Agent} & \textbf{Deep RL Agent} & \textbf{Reasoning Agent} \\
\midrule
Representación interna & Estado simbólico & Embedding neuronal (vector) & Lenguaje natural \\
Forma de construcción & Reglas escritas a mano & Entrenamiento de millones de pasos & Prompting / Fine-tuning \\
Transferencia entre dominios & Muy deficiente & Deficiente & Buena \\
Dimensión de la representación & Conjunto finito de símbolos & Vector de dimensión fija & Espacio lingüístico infinito \\
\bottomrule
\end{tabular}
\caption{Comparación de los tres paradigmas de Agent}
\end{table}

Ventajas del lenguaje natural como representación intermedia:
\begin{itemize}
    \item Aprovecha el conocimiento abundante del preentrenamiento del LLM, sin necesidad de una gran cantidad de entrenamiento adicional
    \item La capacidad expresiva es ilimitada: se puede pensar en una sola palabra o en un millón de tokens
    \item Admite de manera natural el inference-time scaling
\end{itemize}

\subsection{Resumen de la sección}

La innovación esencial del Reasoning Agent consiste en usar el lenguaje natural como representación interna y medio de razonamiento del Agent, lo cual le otorga un espacio de acciones internas de dimensión infinita y la capacidad de transferir el conocimiento preentrenado sin ejemplos.

\section{Memoria a largo plazo (Long-term Memory)}

\subsection{Las limitaciones de la memoria a corto plazo}

La ``memoria a corto plazo'' del Agente LLM es la ventana de contexto (context window), que tiene tres limitaciones fundamentales:

\begin{enumerate}
    \item \textbf{Solo se puede anexar} (append-only): solo se pueden agregar tokens nuevos continuamente, no se pueden editar ni eliminar
    \item \textbf{Capacidad limitada}: incluso con una ventana de un millón de tokens, sigue habiendo un límite
    \item \textbf{No persiste}: al terminar la tarea o al vaciarse el contexto, toda la ``memoria'' desaparece
\end{enumerate}

\begin{warningbox}{El dilema del pez dorado}
Un Agente sin memoria a largo plazo es como un pez dorado: aunque hoy demuestre la conjetura de Riemann, mañana tiene que empezar desde cero. Peor aún, el siguiente intento no necesariamente tendrá éxito.
\end{warningbox}

\subsection{Reflexion: la forma más simple de memoria a largo plazo}

Reflexion es una extensión natural de ReAct:

\begin{enumerate}
    \item El Agente intenta resolver una tarea (por ejemplo, escribir código)
    \item La tarea falla (por ejemplo, no pasa las pruebas unitarias)
    \item El Agente \textbf{reflexiona} sobre la causa del fallo (por ejemplo, ``olvidé manejar el caso límite'')
    \item El resultado de la reflexión se escribe en la memoria a largo plazo
    \item En el siguiente intento se lee la memoria a largo plazo para evitar repetir el error
\end{enumerate}

\begin{knowledgebox}{Un nuevo paradigma de aprendizaje}
Reflexion es, en esencia, un método de \textbf{aprendizaje mediante lenguaje}, distinto del RL tradicional, que actualiza los pesos mediante descenso de gradiente:
\begin{itemize}
    \item No requiere una señal de recompensa escalar, puede aprovechar cualquier retroalimentación textual (errores de compilación, resultados de pruebas, comentarios de un instructor, etc.)
    \item No actualiza parámetros mediante retropropagación, sino que influye en el comportamiento futuro actualizando la memoria textual
\end{itemize}
\end{knowledgebox}

\subsection{Formas más complejas de memoria a largo plazo}

\begin{itemize}
    \item \textbf{Voyager}: aprende habilidades en Minecraft, guarda el código exitoso (por ejemplo, ``fabricar una espada'') en una biblioteca de habilidades que puede invocarse directamente en tareas posteriores.
    \item \textbf{Generative Agents}: 25 Agentes que simulan humanos viven en un pueblo; cada Agente mantiene un registro de eventos (episodic memory) y una memoria semántica (semantic memory) extraída de su experiencia, y estas memorias influyen en su comportamiento posterior.
    \item \textbf{Fine-tuning como memoria a largo plazo}: también se puede hacer fine-tuning directamente sobre los parámetros del modelo, ``escribiendo'' la experiencia en los pesos de la red neuronal.
\end{itemize}

\subsection{KOALA: una abstracción unificada del Agente}

Shunyu Yao (姚顺雨) propuso el marco KOALA, que describe de manera unificada cualquier Agente mediante tres componentes:

\begin{enumerate}
    \item \textbf{Memory}: el lugar donde se almacena la información (ventana de contexto + almacenamiento externo + parámetros del modelo)
    \item \textbf{Action Space}: el conjunto de acciones que el Agente puede ejecutar
    \item \textbf{Decision-making Procedure}: cómo elegir la siguiente acción dados la memoria y el espacio de acciones
\end{enumerate}

\subsection{Resumen de la sección}

La memoria a largo plazo hace que el Agente evolucione de ser un solucionador de un solo uso, ``sin estado'', a un sistema de aprendizaje capaz de acumular experiencia y mejorar de manera continua. La memoria puede ser un registro textual, una biblioteca de código de habilidades o incluso los propios parámetros del modelo.
\section{Nuevas áreas de aplicación}

\subsection{Automatización digital (Digital Automation)}

Los agentes LLM abrieron una dimensión completamente nueva de aplicación: la \textbf{automatización digital}:

\begin{itemize}
    \item \textbf{WebShop}: el agente ejecuta tareas de compra (buscar, comparar, personalizar, comprar) en un sitio de comercio electrónico simulado a partir de Amazon, construido con datos de productos reales a escala de millones.
    \item \textbf{Interacción con páginas web}: el agente opera páginas web reales para completar diversas tareas.
    \item \textbf{SWE-bench}: dado un repositorio de GitHub y un issue, el agente debe entender el código base, escribir pruebas y generar un parche de corrección: una tarea real de ingeniería de software.
\end{itemize}

\subsection{Descubrimiento científico}

Como en el trabajo de CHEM-AI: el agente analiza datos químicos, usa herramientas (Python, motores de búsqueda) y propone compuestos nuevos; las propuestas se validan mediante síntesis en un laboratorio húmedo, con lo que el espacio de acción del agente se extiende al mundo físico.

\subsection{Resumen de la sección}

Las aplicaciones de los agentes LLM se han extendido de las preguntas y respuestas y los juegos propios del NLP/RL tradicional a la automatización digital y el descubrimiento científico reales, algo que los paradigmas de agente anteriores no podían alcanzar.
\section{Direcciones futuras}

Shunyu Yao propuso cinco direcciones de investigación clave:

\subsection{Entrenamiento (Training)}

\begin{importantbox}{Coevolución de modelo y Agent}
Los datos de entrenamiento de los LLM actuales provienen principalmente de texto de internet y carecen de datos de trayectorias de Agent (proceso de razonamiento, secuencia de acciones, retroalimentación del entorno). Solución: usar un prompted agent para generar grandes cantidades de datos de trayectorias, y luego usar esos datos para hacer fine-tuning al modelo, formando un ciclo positivo. Es análogo a la coevolución entre las GPU y el aprendizaje profundo.
\end{importantbox}

\subsection{Interfaz (Interface)}

\begin{itemize}
    \item La interfaz humano-máquina (HCI) se ha estudiado durante décadas, pero la interfaz óptima para un Agent es distinta de la que conviene a los humanos
    \item Por ejemplo: los humanos buscan archivos con \texttt{ls} + \texttt{cd} (un resultado a la vez), pero a un Agent le conviene más obtener varios resultados de una vez
    \item No se trata de optimizar al Agent, sino de optimizar el \textbf{entorno}: el mismo modelo se desempeña mejor con una mejor interfaz
\end{itemize}

\subsection{Robustez (Robustness)}

\begin{warningbox}{Pass@K vs. robustez}
El ámbito académico se enfoca en pass@K (tener éxito al menos una vez en K intentos), pero las aplicaciones reales necesitan \textbf{robustez} (tener éxito en cada uno de los K intentos). En todos los modelos actuales, la robustez disminuye de manera monótona a medida que aumenta el número de muestreos. Un Agent de servicio al cliente puede perder un cliente con un solo fallo: no basta con aspirar a «poder lograrlo», también hay que «poder lograrlo siempre».
\end{warningbox}

\subsection{Colaboración humano-máquina (Human-in-the-loop)}

En escenarios reales, el Agent necesita interactuar con personas: el cliente no da toda la información de una sola vez, y el Agent necesita varias rondas de conversación para obtener la información que requiere.

\subsection{Bancos de pruebas (Benchmarks)}

\begin{itemize}
    \item Se necesitan bancos de pruebas más realistas y de mayor escala (como Tau-Bench: escenarios de servicio al cliente)
    \item Se necesita evaluar la robustez y no solo el límite superior de la capacidad
    \item Se necesita incluir elementos de interacción humano-máquina
\end{itemize}

\subsection{Resumen de la sección}

Las cinco direcciones clave del campo de los Agent de LLM (entrenamiento, interfaz, robustez, colaboración humano-máquina y evaluación) tienen mucha fruta fácil de alcanzar y no requieren cantidades enormes de recursos de GPU, por lo que son adecuadas para la investigación académica.

\section{Síntesis y ampliación}

\subsection{Puntos clave}

\begin{enumerate}
    \item \textbf{Reasoning Agent} es un paradigma de Agente completamente nuevo que usa el lenguaje natural como representación interna y cuenta con un espacio de razonamiento de dimensión infinita.
    \item \textbf{ReAct} unifica el razonamiento y la acción, que se refuerzan mutuamente: el razonamiento guía la planeación de acciones y el manejo de excepciones, mientras que la acción aporta al razonamiento información externa y cómputo.
    \item La \textbf{memoria de largo plazo} (Reflexion, Voyager, Generative Agents) permite que el Agente acumule experiencia entre tareas.
    \item La \textbf{automatización digital} es una dimensión de aplicación completamente nueva que abren los Agentes LLM.
    \item La simplicidad (simplicity) es la cualidad de investigación más poderosa: simple significa general.
\end{enumerate}

\subsection{Lecturas adicionales}

\begin{itemize}
    \item Yao et al., ``ReAct: Synergizing Reasoning and Acting in Language Models,'' ICLR 2023.
    \item Shinn et al., ``Reflexion: Language Agents with Verbal Reinforcement Learning,'' NeurIPS 2023.
    \item Wang et al., ``Voyager: An Open-Ended Embodied Agent with Large Language Models,'' 2023.
    \item Park et al., ``Generative Agents: Interactive Simulacra of Human Behavior,'' UIST 2023.
    \item Yao et al., ``WebShop: Towards Scalable Real-World Web Interaction with Grounded Language Agents,'' NeurIPS 2022.
    \item Zhou et al., ``WebArena: A Realistic Web Environment for Building Autonomous Agents,'' ICLR 2024.
    \item Yao et al., ``Tau-Bench: A Benchmark for Tool-Agent-User Interaction in Real-World Domains,'' 2024.
\end{itemize}

%% --- Fin del contenido --- %%

\end{document}
